\section{Generation locality}\label{sec:gen-locality}

This section shows that the generations of Part~X, built with fixed couplings on tori of growing side, converge at the flat configurations, and that their coefficients converge (Theorem~\ref{thm:gen-locality}). This is what Proposition~\ref{prop:beta-identification} uses. The proof has two parts.
\begin{itemize}
\item For a polymer that does not wrap and at a flat configuration, the new terms of one step are a single functional, defined on the infinite lattice, of the old activities at flat configurations (Proposition~\ref{prop:gl-functional}). This extends the torus independence of the one-loop pieces (Proposition~\ref{prop:one-loop-local}) to the whole step at fixed old data.
\item The old activities converge as the torus grows, by induction over the scales, and the functional is continuous in them. The induction follows the proof of Proposition~\ref{prop:flow-continuity}(b), with the side of the torus in place of the couplings.
\end{itemize}
The generations themselves are not shown to be independent of the torus, and they are not expected to be: the coefficients $\beta_i$ are averages over the polymers that do not wrap \eqref{eq:beta-vol}, and they enter the later generations through the old activities. No rate of convergence is proved (Remark~\ref{rem:gl-rate}).

\paragraph{The result in words, and its place.} Fix couplings $(g_0,\dots,g_{j-1})\in\Sigma_{j-1}$, with $\gamma$ as in ($\Gamma$1)--($\Gamma$3) and $G$ compact, connected and simple, and let the side of the torus grow. At the flat configurations, which equal $e^{\ii\eta B}$ near a polymer and $1$ elsewhere (Section~\ref{subsec:gl-setting}), the raw generation and the generation of every scale $i\le j$ converge to activities of the infinite lattice, the coefficients $\beta_i$ converge to the coefficients $\beta^\infty_i$ of the limits, and the limits depend only on $(g_0,\dots,g_{i-1})$ (Theorem~\ref{thm:gen-locality}). The reason is that every construction of a step is given by local rules: on a polymer that does not wrap, the new terms are a functional of local data, which are the same on all large tori (Proposition~\ref{prop:gl-functional}). Theorem~\ref{thm:gen-locality} is stated at the end of the section, after the transfer lemmas. This section is off the main line: Theorem~\ref{thm:ol-uv}, and so Theorem~\ref{thm:M}(iv), does not use it. Its only use is Proposition~\ref{prop:beta-identification}, the comparison of the $\beta_j$ with the coefficients that Ba\l aban's definition assigns to the generations.

\subsection{The setting}\label{subsec:gl-setting}

\paragraph{Setting.} Let $G$ be compact, connected and simple. The parameters are those of Section~\ref{sec:flow4} (Setting), and $\gamma$ satisfies ($\Gamma$1)--($\Gamma$3) of Remark~\ref{rem:flow-constants}. Let $j\ge1$ and $\pi=(g_0,\dots,g_{j-1})\in\Sigma_{j-1}$. For $i\le j$ and $m\in\mathfrak N_i$ let $E'^{(i)}[m]$, $E^{(i)}[m]=E'^{(i)}[m]^{\rm sym}$ and $\beta_i[m]=\beta_{E^{(i)}[m]}$ be the raw generation, the generation and the coefficient of scale $i$ constructed with $(g_0,\dots,g_{i-1})$ on the torus whose lattice of scale $i$ has side $m$ (Proposition~\ref{prop:flow-continuity}, with $K=i$).

This is legitimate. The Setting of Section~\ref{sec:generations}, which defines $\mathfrak N_i$, includes that of Section~\ref{subsec:tf-scales}, and so $n_K\ge n_{\min}$ with $K=i$. So the Setting of Section~\ref{sec:flow4} holds with $K=i$, and Proposition~\ref{prop:flow-continuity}(a) applies. By that proposition, $E^{(l)}$ and $\beta_l$, $l\le i$, depend only on $(g_0,\dots,g_{l-1})$ and on the torus; the steps $0,\dots,l-1$ do not use the scales above $l$. By the proof of Theorem~\ref{thm:ol-cumulative}, the side of this torus at scale $l$ is $L^{i-l}m\in\mathfrak N_l$. So the generations of scale $l$ on the torus of side $m$ at scale $i$ are the $E^{(l)}[L^{i-l}m]$. For $1\le k<j$ we write $\mathcal E_k[m']$ for the old activities of step $k$ (Definition~\ref{def:old-activities}) on the torus whose lattice of scale $k+1$ has side $m'\in\mathfrak N_{k+1}$. They are built from the $E^{(l)}[L^{k+1-l}m']$ and the $\beta_l[L^{k+1-l}m']$, $l\le k$.

\paragraph{Flat configurations.} Use the units of scale $i$. For a polymer $X$ of scale $i$ of the infinite lattice, or of a torus on which it does not wrap, $\mathcal U^\flat_1(X)$ is the flat domain of Section~\ref{subsec:ol-period-vacpol}. Let $\mathcal U_1(X)$ be the set defined in the same way with $\delta_1$ in place of $\delta_1^\flat$. It consists of the configurations $V_B$ equal to $e^{\ii\eta B}$ on the bonds of $X^{(3)}$, read in a lift, and to $1$ elsewhere, with $|B|<\delta_0$ and $|\nabla B|+[\nabla B]_\alpha<\delta_1$ on $X^{(3)}$; the differences and quotients are taken with the real part $1$. By Definition~\ref{def:graded} with the real part $1$,
\[
\mathcal U^\flat_1(X)\subset\mathcal U_1(X)\subset\mathcal U(X).
\]
Their points are called \emph{flat}. As in Section~\ref{subsec:ol-period-vacpol}, a flat point is identified with its field $B$, and on the torus whose lattice of scale $i$ has side $m$ we write $V_B[m]$. The conditions are imposed on the bonds of $X^{(3)}$ only. So they are the same on the infinite lattice and on every torus on which $X$ does not wrap. As sets of fields $B$, $\mathcal U_1(X)$ and $\mathcal U^\flat_1(X)$ are open and convex in a finite-dimensional space, since the bounds are strict and are bounds on seminorms of $B$.

\emph{The infinite lattice.} Polymers of scale $i$ of the infinite lattice, their tree lengths and their numbers of cubes are those of Section~\ref{subsec:ol-period-vacpol}. Their closures $\operatorname{cl}_k$ at the scales $k\ge i$, the sets $\Box^+$, the base cubes and the classes of Definition~\ref{def:polymer-classes} are defined by the same rules as on the torus. The base cube of a polymer of the infinite lattice is read in the box spanned by its cubes. These rules are local in the grids. So for a polymer that does not wrap they commute with the projection to the torus (proof of Lemma~\ref{lem:OL}; Section~\ref{sec:flow}, base cubes).

\begin{lemma}[Flat points]\label{lem:gl-flat}
\begin{enumerate}[label=(\alph*)]
\item Let $l\le k$, let $Y$ be a polymer of scale $k$ and $X\subset Y$ a polymer of scale $l$, both not wrapping or both of the infinite lattice. A point of $\mathcal U^\flat_1(Y)$ of scale $k$ is, on $X^{(3)}$ and in the units of scale $l$, a point of $\mathcal U^\flat_1(X)$ of scale $l$. Let $Y'$ be a polymer of scale $k+1$, not wrapping or of the infinite lattice. A point of $\mathcal U_1(Y')$ of scale $k+1$ lies, on $X^{(3)}$ for every polymer $X\subset Y'$ of scale $k$, in $\mathcal U^\flat_1(X)$ of scale $k$, with budgets $L^{-2}\delta_1$ and $L^{-1}\delta_0$.
\item Let $Y'$ be a polymer of scale $k+1$ that does not wrap, and $V=V_B\in\mathcal U_1(Y')$. With the representation $(1,B)$, the cutoff extension of Remark~\ref{rem:bg-to-graded} is $U'=e^{\ii\eta\chi B}$. It has real part $1$, it equals $V$ on the bonds within unit distance of $Y'$, in the units of scale $k+1$, and it equals $1$ off $Y'^{(3)}$. Read in the units of step $k$, it is a reference of the Setting of Section~\ref{sec:step-quantities}.
\item Let $r\in\mathfrak S^{\rm grid}_i$, and let $X$ be a polymer of scale $i$ that does not wrap. Then $r\,V_B=V_{rB}$, with $(rB)(b)=B(r^{-1}b)$ and $B(-b)=-B(b)$, and $r\,\mathcal U^\flat_1(X)\subset\mathcal U_1(rX)$.
\end{enumerate}
\end{lemma}

\begin{proof}\hasdetail{lem:gl-flat}
(a) This is Remark~\ref{rem:domain-constants}(i), with $\delta_1^\flat$ in place of $\delta_1$ for the first statement. Its proof keeps the real part, here $1$.

(b) This is Remark~\ref{rem:bg-to-graded} with $U_r=1$. The cutoff $\chi$ equals $1$ at unit distance at most $1$ from $Y'$ and $0$ at distance at least $2$. The last statement is Remark~\ref{rem:domain-constants}(ii): the real part $1$ is regular, with curvature and current $0$.

(c) Apply Lemma~\ref{lem:grid-domains} with $U_r=1$, $\delta'_1=\delta_1^\flat$ and $\delta'_0=\delta_0$. Then $R(U_r(b))=1$, so $rA$ is as stated, and $rU_r=1$. By (DC6), $2\delta_1^\flat+C_{\rm refl}\theta_r\delta_0\le\delta_1$, as in the derivation of \eqref{eq:flat-in-full}. So $rV_B=e^{\ii\eta\,rB}$ on the bonds of $(rX)^{(3)}$, with the budgets of $\mathcal U_1(rX)$, and $rV_B=1$ on the other bonds.
\end{proof}

\paragraph{Admissible sides.} Let $k\ge0$, and let $\tilde Y'$ be a polymer of scale $k+1$ of the infinite lattice. A side $n$ at step $k$ is \emph{admissible for $\tilde Y'$} if $n/L\in\mathfrak N_{k+1}$ and $\operatorname{diam}\tilde Y'\le\frac12(n/L)-8$ in the units of scale $k+1$. Then the projection $Y'=\pi_n\tilde Y'$ to the torus $\Tf[n]$ of step $k$ does not wrap at scale $k+1$. By Lemma~\ref{lem:torus}(a), $n$ is admissible at step $k$ in the sense of Section~\ref{subsec:torus-independence}, as in the proof of Lemma~\ref{lem:OL}. In the units of step $k$ the lift $\tilde Y'$ has diameter at most $\frac12n-8L$. So, for two sides $n$, $n'$ admissible for $\tilde Y'$ and every $a\in M\Z^4$, the transfer of the lift $\tilde Y'$ with the vector $a$ is a grid transfer from $\Tf[n]$ to $\Tf[n']$ (Definition~\ref{def:grid-transfer}(ii)), with $P=Y'$; here $n'=n$ is allowed. Every $\tilde Y'$ has admissible sides, since $\mathfrak N_{k+1}$ is unbounded.

\subsection{Transfer of one step}\label{subsec:gl-transfer}

This subsection uses the setting of Section~\ref{subsec:torus-independence} at step $k$, with its grid transfers and the consequence of Lemma~\ref{lem:grid-transfer}. Lemma~\ref{lem:local-inventory}(c) lists the linear operators of the step, the elimination $\ell_c$, $\zeta_c$, $\hat Z^1$, $\hat Z^2$ (hence the map $\hat Z$ of Lemma~\ref{lem:eliminate}(d)), the transport $\mathfrak G_W$, the curvatures $F_U$, $F^s_U$ and the current $J_U$, and the averaged backgrounds $W_j$, at complex references through Proposition~\ref{prop:AB}(i). The step also uses the following nonlinear maps.

\begin{lemma}[Nonlinear local maps]\label{lem:gl-inventory}
Each of the following is given on every admissible torus by the same rule. At a site based at $y$ the rule uses only the objects of Lemma~\ref{lem:local-inventory}(a) and the configuration and the fields on the bonds within distance $4L+2\le R_{\rm g}$ of $y$. At complex references it is the holomorphic continuation of the same formula.
\begin{enumerate}[label=(\roman*)]
\item The nonlinear part $\Cav_k$ of the gauge-adapted averaging. The value $\Cav_k(A)(c)$ is a function of $A$ and of the reference on the bonds of $S_c=B(c_-)\cup B(c_+)$.
\item The plaquette terms $s_p$ of the action, and the plaquette terms $n_p$ of $\mathcal N$ (Proposition~\ref{prop:action}). Each is a function of the configuration, respectively of $A$ and of the reference, on the bonds of $p$. Hence $\nabla\mathcal N(A)(b)$ is a function of $A$ and of the reference on the plaquettes that contain $b$.
\item The functions $H_{y,x}$ and $f\circ H_{y,x}$ of Definition~\ref{def:gf-pieces}, functions of $Z$ and of the base on the inner bonds of $B'(y)$.
\item The terms $\log(\sigma(Z(b))/\sigma(0))$ and $\log(\mathsf J_c(Z)/\mathsf J_c(0))$ of $\mathfrak J$ (Proposition~\ref{prop:jacobians}(a)), functions of $Z$ and of the base within distance $2L$ of one bond.
\item The bondwise operations $\ell$ and $\oplus$, the latter the Baker--Campbell--Hausdorff combination of the proof of Proposition~\ref{prop:composed}.
\end{enumerate}
Consequently, let $\iota:P\to P'$ be a grid transfer, and let the references and the fields agree under $\iota$ on $P$. Then each of these objects has the same value at a site $x$ based in $P$ whose range lies in $P$ and at $\iota x$.
\end{lemma}

\begin{proof}\hasdetail{lem:gl-inventory}
(i) By Proposition~\ref{prop:Q-expansion}, $\Cav_k(A)=\delta_k^{-1}\ell(\mathcal Q_k(e^{\ii\eta A}))-\bQ_kA$, with the gauge-adapted averaging $\mathcal Q_k$ of Section~\ref{sec:adapted}. It is defined by the same formulas on every torus, through the relative one-step averages at the backgrounds $W_j$ of the reference. By Proposition~\ref{prop:Q-expansion}(a), Proposition~\ref{prop:linearized-averages}(d) and Lemma~\ref{lem:one-step-basic}, as in the proof of Theorem~\ref{thm:weak-bg}(d), $\Cav_k(A)(c)$ depends only on $A$ and on the reference on the bonds of $S_c$. At complex references, $W_j(U)=\mathfrak a_j(W_{j-1}(U))$ is the same one-step recursion (proof of Lemma~\ref{lem:AB-rep}(b)), and Proposition~\ref{prop:AB}(iii) gives Proposition~\ref{prop:Q-expansion} there.

(ii) The term $s_p$ is explicit in $U(\partial p)$ (Section~\ref{sec:action}). The plaquette terms of $\mathcal N=\sum_pn_p$ are functions of $A$ and of $U_0$ on the bonds of $p$ (proof of Theorem~\ref{thm:weak-bg}(d)). At complex references both are the continuations of Proposition~\ref{prop:complex-ext}(c).

(iii) This is Proposition~\ref{prop:gf-pieces}(d), with the averaged contour variables of the box spanned by $y$ and $x$, which are defined at complex bases by (DC5) (Remark~\ref{rem:domain-constants}(ii)).

(iv) This is Proposition~\ref{prop:jacobians}(a). The function $\sigma$ is explicit, and
\[
\mathsf J_c(Z)=\frac{\sigma(\zeta_c(Z))}{\sigma(0)\,\det\partial\ell_c/\partial Z(b_0(c))}\qquad\text{at }Z(b_0(c))=\zeta_c(Z)
\]
(Proposition~\ref{prop:exact-step}) is built from the functions $\ell_c$ and $\zeta_c$ of Lemma~\ref{lem:local-inventory}(c).

(v) These are bondwise.

The ranges are at most $4L+2$, as in the proof of Lemma~\ref{lem:local-inventory}(c). The consequence follows from Lemma~\ref{lem:grid-transfer}.
\end{proof}

\begin{lemma}[Transfer of the weakened background map]\label{lem:gl-bg}
Let $\iota:P\to P'$ be a grid transfer from $\Tf[n]$ to $\Tf[n']$ at step $k$, and let $U_0$, $U'_0$ be references as in Section~\ref{sec:weak-background}, real or complex as in Proposition~\ref{prop:complex-ext}(d), that agree on $P$. Let $Z$ be a union of boxes with $Z^{(r+4)}\subset P$, let $s\in\mathbb D_\kappa$ with $s_\Box=0$ for $\Box\not\subset Z$, and let $s'$ be given by $s'_{\iota\Box}=s_\Box$ for $\Box\subset Z$ and $s'_{\Box'}=0$ for the other boxes $\Box'$. Let $B$, $B'$ be coarse fields with $|B|_\infty,|B'|_\infty\le R_P$ that agree under $\iota$ on the coarse bonds of $Z^{(r+1)}$. Then $A_1(s';B')$ and $A_*(s';B')$, at $U'_0$, vanish off $\iota(Z)^{(r)}$, and on its bonds they are the transfers by $\iota$ of $A_1(s;B)$ and $A_*(s;B)$, at $U_0$.
\end{lemma}

\begin{proof}\hasdetail{lem:gl-bg}
Write $\iota_*$ for the transfer of fields and operators (Definition~\ref{def:grid-transfer}). We follow the proof of Theorem~\ref{thm:weak-bg}(d) on both tori. Its class $\mathcal K$ is defined on each torus with $Z$, respectively $\iota(Z)$.

\emph{The operators.} The operators $\bC(s)$, $\bH(s)$, $D\bH(s)$ and their transposes are sums of products, as in Lemma~\ref{lem:weak-chain}, of the families of the levels (L1)--(L4) on the domain $\Tf$, with $W=0$ at (L3), and of the local operators $D$, $D^T$, $\bQ$, $\bQ^T$ (Section~\ref{sec:tiers}, the definitions of $\bH(s)$ and $\bC(s)$ on the torus). By Lemma~\ref{lem:weak-basic}(b), at $s$ only their terms with $X_\omega\subset Z$ contribute, and at $s'$ only the images of these. Since $X_\omega^{(r)}\subset Z^{(r)}\subset P$, Lemma~\ref{lem:torus-transfer}(b) carries the factors, and Step~4 of its proof carries the products: $\iota(\omega)$ is a term with polymer $\iota(X_\omega)$ and the carried monomial, and $O'_{\iota(\omega)}=\iota_*O_\omega$. The local operators are carried by Lemma~\ref{lem:local-inventory}(c). Hence $\bC'(s')=\iota_*\bC(s)$, and likewise for $\bH$, $D\bH$ and the transposes, on the fields supported in $Z^{(r)}$ and the coarse fields on the coarse bonds of $Z^{(r)}$.

\emph{The linearizing map.} Let $A'$ be a fine field that vanishes off $Z^{(r)}$, and let $A''$ vanish off $\iota(Z)^{(r)}$ and agree with $\iota_*A'$ there. The iterates of \eqref{eq:Ds} from $0$, $X_{p+1}=\Cav_k(A'-\bH(s)X_p)$, stay in $\mathcal K$, so they vanish off $Z^{(r+1)}$. By induction on $p$, $X''_p=\iota_*X_p$. The step uses the operators above and Lemma~\ref{lem:gl-inventory}(i) at the coarse bonds starting in $Z^{(r+1)}$, whose sets $S_c$ lie in $Z^{(r+3)}\subset P$. The iterates converge to $\mathcal D_s(A')$, since \eqref{eq:Ds} is a contraction of a ball (proof of Lemma~\ref{lem:Ds}). So $\mathcal D'_{s'}(A'')=\iota_*\mathcal D_s(A')$. The derivative $d\mathcal D_s^T=d\Cav_k(A)^T(I+\bH(s)^Td\Cav_k(A)^T)^{-1}$ is carried term by term in its Neumann series (proof of Lemma~\ref{lem:Ds}). Its terms are products of carried operators, since $d\Cav_k(A)$ is the derivative of a carried local map.

\emph{The fixed point.} The iterates $A_{1,p+1}=\mathcal T_s(A_{1,p})$, $A_{1,0}=0$, stay in $\mathcal K$. By \eqref{eq:Ts} and \eqref{eq:gradV-s}, $\mathcal T_s(A_1)=-\bC(s)\nabla\mathcal V_s(A_1+\bH(s)B)$. The gradient is built from $\mathcal D_s$ and $d\mathcal D_s^T$ at $A_1+\bH(s)B$, the operators above, $\nabla\mathcal N$ (Lemma~\ref{lem:gl-inventory}(ii)) and $\mathsf F=F^s_{U_0}$ (Lemma~\ref{lem:local-inventory}(c)). The field $\bH(s)B$ depends on $B$ only on the coarse bonds of $Z^{(r)}$, where $B$ and $B'$ agree. So, by induction on $p$, $A''_{1,p}=\iota_*A_{1,p}$. The iterates converge in $|\cdot|_*$, hence bondwise, to the fixed point $A_1(s;B)$, since $\mathcal T_s$ is a contraction of the ball of Theorem~\ref{thm:weak-bg}(a). Then $A_*(s;B)=A'_s-\bH(s)\mathcal D_s(A'_s)$, $A'_s=A_1(s;B)+\bH(s)B$, is carried. At complex references the operators are those of Proposition~\ref{prop:complex-ext}(d), and the class $\mathcal K$ is preserved by the same proof (proof of Proposition~\ref{prop:complex-ext}, (viii)).
\end{proof}

\begin{lemma}[Composed activities at flat references]\label{lem:gl-composed}
Let $k\ge1$, and let $X$ be an old polymer of step $k$ that does not wrap. Let $\mathcal E_X$ be a holomorphic function on $\mathcal U^\flat_1(X)$ of step $k$, as a function of the field on the bonds of $X^{(3)}$, which depends only on the bonds of $X$ and is bounded by $\|\mathcal E_X\|$. Let $U'=e^{\ii\eta A'}$ be a reference of the Setting of Section~\ref{sec:step-quantities} with real part $1$, such that the field $A'$ has, on $X^{(3)}$, the budgets $\delta'_1<\delta_1^\flat$ and $\delta'_0<\delta_0$. Let $0<\zeta\le\frac12R_P$, let $\tau$ be as in Proposition~\ref{prop:composed} with $\delta_1^\flat$ in place of $\delta_1$, and assume $\tau<\frac12$. Then for $s\in\mathbb D_\kappa$ and $|\hat Z(Z)|_\infty\le\zeta$:
\begin{enumerate}[label=(\alph*)]
\item on $X^{(3)}$, the configuration $e^{\ii\eta A_*(s;\hat Z(Z))}U'$ is a point of $\mathcal U^\flat_1(X)$. The function
\[
f(X;s,Z)=\mathcal E_X\big(e^{\ii\eta A_*(s;\hat Z(Z))}U'\big)-\mathcal E_X(U')
\]
is analytic in $(s,Z)$ and in $A'$ jointly, continuous up to the boundary in $s$, and $|f(X;s,Z)|\le4\tau\|\mathcal E_X\|$;
\item if $s_\Box=0$ outside a union of boxes $Z'$ that contains the boxes meeting $X$, then $f(X;s,Z)$ depends on $s$ only through the boxes of the $*$-component $Z'_0$ of $Z'$ containing them, and on $Z$ and $U'$ only on $(Z'_0)^{(r')}$.
\end{enumerate}
\end{lemma}

\begin{proof}\hasdetail{lem:gl-composed}
The proof of Proposition~\ref{prop:composed} applies with $P=X$, $r_E=0$ and $U'_r=1$. It evaluates $\mathcal E_X$ only at the configurations $e^{\ii\eta tA_*}U'=e^{\ii\eta(tA_*\oplus A')}$ on $X^{(3)}$, $|t|<\frac1{2\tau}$, whose real part is $U'_r=1$, and it shows that they satisfy the bounds of $\mathcal U(X;\theta_r,\delta_1^\flat,\delta_0)$ strictly, with the representation $(1,tA_*\oplus A')$. These are points of $\mathcal U^\flat_1(X)$. Lemma~\ref{lem:norm-conversion} is applied with $U'_r=1$.
\end{proof}

\begin{lemma}[Transfer of the pieces]\label{lem:gl-pieces}
Let $k\ge0$, and let $\iota:P\to P'$ be a grid transfer from $\Tf[n]$ to $\Tf[n']$ at step $k$. Let $U'$ and $U''$ be references of the Setting of Section~\ref{sec:step-quantities} on the two tori that agree on $P$, at the same coupling $g\in(0,\gamma]$. Let $Z$ be a polymer of step $k$ with $Z^{(r')}\subset P$, and let $\xi$, $\xi'$ be white-noise fields with $|\xi|_\infty,|\xi'|_\infty\le p_0(g)$ and $\xi'(\iota b)=\xi(b)$ for the bonds $b\in\mathfrak F$ in $Z^{(r')}$.
\begin{enumerate}[label=(\alph*)]
\item The pure pieces of Proposition~\ref{prop:localized-pieces} satisfy $E'_{\rm p}(\iota Z)(\xi')=E_{\rm p}(Z)(\xi)$. At $k=0$ this holds for the pure pieces of the first step (Proposition~\ref{prop:first-step}).
\item Let $k\ge1$, and let $U'$ and $U''$ have real part $1$. For every old polymer $X\subset Z$ let $\mathcal E_X$ and $\mathcal E'_{\iota X}$ be functions as in Lemma~\ref{lem:gl-composed}, with $\mathcal E'_{\iota X}(\iota_*A)=\mathcal E_X(A)$ for the flat fields $A$ on the bonds of $X^{(3)}$, and let $U'$ satisfy the hypothesis of Lemma~\ref{lem:gl-composed} on $X^{(3)}$. Let $E_{\rm m}(Z)=\sum_{X\subset Z}f_{\rm loc}(X,Z)$ be built from the functions $f(X;s,Z(s))$ of Lemma~\ref{lem:gl-composed} as in Proposition~\ref{prop:localized-pieces}, and $E'_{\rm m}(\iota Z)$ likewise. Then $E'_{\rm m}(\iota Z)(\xi')=E_{\rm m}(Z)(\xi)$.
\end{enumerate}
\end{lemma}

\begin{proof}\hasdetail{lem:gl-pieces}
(a) Let $\Box\subset Z$, and let $f$ be one of the functions localized in the proof of Proposition~\ref{prop:localized-pieces} with $Y=\Box$: $V(\Box;s)$, $G_3(\Box;s)$, or the terms of $\mathfrak J(Z(s))$ grouped in $\Box$. Let $f'$ be the same function on $\Tf[n']$ for the box $\iota\Box$. By Lemma~\ref{lem:localization}, $f_{\rm loc}(\Box,Z)=\int_{[0,1]^{Z\setminus\Box}}\partial_{s_{Z\setminus\Box}}f(s(Z,t))\,dt$, and the points $s(Z,t)$ vanish off $Z$. We show $f'(s')=f(s)$ for every $s\in\mathbb D_{\kappa'}$ that vanishes off $Z$, with $s'$ as in Lemma~\ref{lem:gl-bg}. Then $f$ and $f'$ are the same analytic functions of the variables $s_\Box$, $\Box\subset Z$, so $f'_{\rm loc}(\iota\Box,\iota Z)=f_{\rm loc}(\Box,Z)$, and summing over $\Box\subset Z$ gives (a).

Let $s$ vanish off $Z$. By Lemma~\ref{lem:weak-basic}(b), only the terms of $C_k^{1/2}(s)$ with $X_\omega\subset Z$ contribute, and the operator vanishes off $Z^{(r)}\times Z^{(r)}$. Lemma~\ref{lem:torus-transfer}(b) carries these terms (the family $C_k^{1/2}(s)$ on $\Tf$, Corollary~\ref{cor:weak-covariance}), since $X_\omega^{(r)}\subset Z^{(r)}\subset P$. So $Z'(s')=gC'^{1/2}_k(s')\xi'$ vanishes off $\iota(Z)^{(r)}$, $Z(s)$ vanishes off $Z^{(r)}$, and $Z'(s')(\iota b)=Z(s)(b)$ for the bonds $b$ of $P$.

\emph{Interaction.} $V(\Box;s)=g^{-2}\sum_pv_p(s;Z(s))$, the sum over the plaquettes $p$ based in $\Box$. By Proposition~\ref{prop:interaction-pieces}(d), $v_p=0$ unless $p$ has a point in $Z^{(r+1)}$, and then $p$ lies in $Z^{(r+2)}$. For fields $Y$ on $\mathfrak F$ that vanish off $Z^{(r)}$, with $|Y|_\infty<r_I$, put $h(Y)=s_p(e^{\ii\eta A_*(s;\hat Z(Y))}U')$.
\begin{itemize}
\item The field $\hat Z(Y)$ equals $Y$ on $\mathfrak F$ and $\zeta_c(Y)$ on $b_0(c)$. Here $\zeta_c(Y)$ is a function of $Y$ and of the base on the bonds of $B'(c_-)\cup B'(c_+)$ (Lemma~\ref{lem:eliminate}(c)), carried by Lemma~\ref{lem:local-inventory}(c). So $\hat Z'(\iota_*Y)$ and $\hat Z(Y)$ agree under $\iota$ on the coarse bonds of $Z^{(r+1)}$. Moreover $|\hat Z(Y)|_\infty<R_P$, by the choice of $r_I$ (Section~\ref{subsec:sq-interaction}).
\item By Lemma~\ref{lem:gl-bg}, applicable since $Z^{(r+4)}\subset Z^{(r')}\subset P$ (as $r'=r+2L+2\ge r+8$), $A_*(s';\hat Z'(\iota_*Y))$ is the transfer of $A_*(s;\hat Z(Y))$.
\item By Lemma~\ref{lem:gl-inventory}(ii), $h'(\iota_*Y)=h(Y)$.
\end{itemize}
The Taylor polynomial $T_2$ at $Y=0$, evaluated at a field that vanishes off $Z^{(r)}$, uses only the derivatives of $h$ at $0$ in the directions of these fields. So $v'_{\iota p}(s';\iota_*Y)=v_p(s;Y)$, and with $Y=Z(s)$, $V'(\iota\Box;s')=V(\Box;s)$.

\emph{Gauge fixing.} $G_3(\Box;s)=g^{-2}\sum_{y\in\T'\cap\Box}\sum_{x\in B'(y)}\varphi_{y,x}(Z(s))$. The function $\varphi_{y,x}=f\circ H_{y,x}-\frac12\ip{\mathfrak G_W\cdot}{\mathfrak G_W\cdot}$ is carried by Lemma~\ref{lem:gl-inventory}(iii) and Lemma~\ref{lem:local-inventory}(c). The blocks $B'(y)$, $y\in\Box$, lie within distance $L$ of $\Box$.

\emph{Jacobians.} The terms of $\mathfrak J(Z(s))$ grouped in $\Box$ are carried by Lemma~\ref{lem:gl-inventory}(iv).

\emph{The first step.} At $k=0$ there is no background map. The interaction pieces are $v_p(Z)=s_p(e^{\ii\hat Z(Z)}U')-T_2[\cdots](Z)$, local within range $2L+2$, and $C_0^{1/2}(s)$ is given by the walk of Theorem~\ref{thm:weak-inverse} with $\fA=\Delta_0+r$ and $\mathfrak N=0$ (Proposition~\ref{prop:first-step}(a)--(c)). Lemma~\ref{lem:torus-transfer} carries it in its case $k=0$. The same argument applies, with the factorization of the composed pieces in the proof of Proposition~\ref{prop:first-step} in place of Proposition~\ref{prop:interaction-pieces}(d).

(b) Let $X\subset Z$ be an old polymer. In the proof of Proposition~\ref{prop:localized-pieces}, $f(X;s,Z(s))$ is localized by Lemma~\ref{lem:localization} with $Y$ the set of boxes of $X$. Let $s$ vanish off $Z$. By Lemma~\ref{lem:gl-composed}(b), $f(X;s,Z(s))$ depends on $Z(s)$ and on $U'$ only on $Z^{(r')}\subset P$. By the first part of the proof, $Z(s)$ is carried. By Lemma~\ref{lem:gl-bg}, so is $A_*(s;\hat Z(Z(s)))$, and the operation $\oplus$ is bondwise. Since $X^{(3)}\subset Z^{(r')}$, the composed flat field on $(\iota X)^{(3)}$ is the transfer of that on $X^{(3)}$. By the hypothesis on $\mathcal E_X$, $f'(\iota X;s',Z'(s'))=f(X;s,Z(s))$. As in (a), $f'_{\rm loc}(\iota X,\iota Z)=f_{\rm loc}(X,Z)$. The map $X\mapsto\iota X$ is a bijection of the old polymers in $Z$ onto those in $\iota Z$, and (b) follows.
\end{proof}

\subsection{The step on the infinite lattice}\label{subsec:gl-functional}

\begin{definition}[Flat old data]\label{def:gl-flat-data}
Let $k\ge1$, and let $\tilde Y'$ be a polymer of scale $k+1$ of the infinite lattice, read in the units of step $k$. \emph{Flat old data on $\tilde Y'$} is a family $\mathcal E^{\rm fl}=(\mathcal E^{\rm fl}(\tilde Y,\cdot))$, indexed by the polymers $\tilde Y\subset\tilde Y'$ of scale $k$ of the infinite lattice. Each $\mathcal E^{\rm fl}(\tilde Y,\cdot)$ is a holomorphic function on $\mathcal U^\flat_1(\tilde Y)$ of scale $k$, depends only on the field on the bonds of $\tilde Y$, vanishes at $0$, and satisfies
\[
\sup_{\mathcal U^\flat_1(\tilde Y)}|\mathcal E^{\rm fl}(\tilde Y,\cdot)|\le C_{\rm old}\mathsf K\,e^{-\kappa_nd_k(\tilde Y)},
\]
with $C_{\rm old}$ of Proposition~\ref{prop:old-activities}. We write $\mathfrak D_k(\tilde Y')$ for the set of such families. For $a\in LM\Z^4$ in the units of step $k$, that is, a vector of $M\Z^4$ in the units of scale $k+1$, $(\tau_a\mathcal E^{\rm fl})(\tilde Y+a,\tau_aB)=\mathcal E^{\rm fl}(\tilde Y,B)$, with $(\tau_aB)(b)=B(b-a)$, defines $\tau_a\mathcal E^{\rm fl}\in\mathfrak D_k(\tilde Y'+a)$.

Let $n$ be a side admissible for $\tilde Y'$, and let $\mathcal E$ be old activities of step $k$ on $\Tf[n]$ as in Lemma~\ref{lem:raw-new}, with $\|\mathcal E\|\le C_{\rm old}\mathsf K$. Their \emph{flat restriction} to $\tilde Y'$, also written $\mathcal E^{\rm fl}_{\tilde Y'}[\mathcal E]$, is
\[
\mathcal E^{\rm fl}_{\tilde Y'}(\tilde Y,B)=\mathcal E(\pi_n\tilde Y,V_B[n]),\qquad\tilde Y\subset\tilde Y',\ B\in\mathcal U^\flat_1(\tilde Y).
\]
\end{definition}

Since $\pi_n\tilde Y'$ does not wrap, neither does $\pi_n\tilde Y$. The configuration $V_B[n]$ lies in $\mathcal U^\flat_1(\pi_n\tilde Y)\subset\mathcal U^\flat(\pi_n\tilde Y)$. The old activities are normalized graded activities at their polymers on these domains, with the stated bound. Moreover $d_k(\pi_n\tilde Y)=d_k(\tilde Y)$, since $\operatorname{diam}\tilde Y\le\frac12n-8$; this is the first part of the proof of Lemma~\ref{lem:gl-limits}(a), which uses only the geometry of polymers. So $\mathcal E^{\rm fl}_{\tilde Y'}\in\mathfrak D_k(\tilde Y')$.

\begin{definition}[The step on the infinite lattice]\label{def:gl-functional}
Let $k\ge0$ and $g\in(0,\gamma]$. Let $\tilde Y'$ be a polymer of scale $k+1$ of the infinite lattice, let $\mathcal E^{\rm fl}\in\mathfrak D_k(\tilde Y')$ if $k\ge1$, and let $B\in\mathcal U_1(\tilde Y')$, of scale $k+1$. Choose a side $n$ admissible for $\tilde Y'$, put $Y'=\pi_n\tilde Y'$, and let $U'$ be the cutoff extension of $V_B[n]$ with the representation $(1,B)$ (Lemma~\ref{lem:gl-flat}(b)), read in the units of step $k$. Consider the pieces
\[
E(Z)=E_{\rm p}(Z)+E_{\rm m}(Z),\qquad E_{\rm m}(Z)=\sum_{X\subset Z}f_{\rm loc}(X,Z),
\]
indexed by the polymers $Z$ of step $k$ with $Z^{(r')}\subset Y'$, at the reference $U'$ and the coupling $g$, and put $E(Z)=0$ for the other polymers. Here:
\begin{itemize}
\item $E_{\rm p}(Z)$ is the pure piece of Proposition~\ref{prop:localized-pieces};
\item the sum runs over the old polymers $X\subset Z$;
\item $f_{\rm loc}(X,Z)$ is given by Lemma~\ref{lem:localization} on $\mathbb D_{\kappa'}$, with $Y$ the set of boxes of $X$, applied to $s\mapsto f(X;s,Z(s))$ of Lemma~\ref{lem:gl-composed}, with $\mathcal E_X(V)=\mathcal E^{\rm fl}(\tilde X,\tilde B_V)$, where $\tilde X$ is the lift of $X$ in $\tilde Y'$ and $\tilde B_V$ is the lift of the flat field of $V$ on $X^{(3)}$.
\end{itemize}
At $k=0$ there are only the pure pieces of the first step (Proposition~\ref{prop:first-step}), with the constants of the first step. Let $E^\#(Y;U')$ be the outputs of Theorem~\ref{thm:mayer} for this system, with $\kappa=\kappa_n-1$ and $\mathsf d=d_k$, and let $E^\#(Y;1)$ be the outputs of the same construction at $B=0$, where $U'=1$. Put
\[
\begin{aligned}
\mathscr S_{k,\tilde Y'}(g;\mathcal E^{\rm fl};B)={}&\sum_{\tilde X:\,\operatorname{cl}_{k+1}\tilde X=\tilde Y'}E^\infty_k(\tilde X,B_{\tilde X})+\sum_{\tilde c:\,\tilde Y'_{\tilde c}=\tilde Y'}J^\infty_k(\tilde c,B_{\tilde c})\\
&+\sum_{Y:\,\hat Y=Y'}\big(E^\#(Y;U')-E^\#(Y;1)\big),
\end{aligned}
\]
where $E^\infty_k$ and $J^\infty_k$ are the infinite-volume pieces of Proposition~\ref{prop:one-loop-local}(d), at the restrictions $B_{\tilde X}$, $B_{\tilde c}$ of $B$, read in the units of step $k$. The first two sums are those of the proof of Lemma~\ref{lem:OL}. At $k=0$ we write $\mathscr S_{0,\tilde Y'}(g;B)$.
\end{definition}

The construction is defined.
\begin{itemize}
\item By Lemma~\ref{lem:gl-flat}(a),(b), $U'$ is a reference of the Setting of Section~\ref{sec:step-quantities}. On $X^{(3)}$ its field has budgets $L^{-2}\delta_1$ and $L^{-1}\delta_0$, for every old polymer $X\subset Y'$. As in the paragraph after Proposition~\ref{prop:composed}, the slack is at least $\frac18\delta_1$. So Lemma~\ref{lem:gl-composed} applies with $\zeta=w$ and $\tau\le C_PL^dgp_0(g)/\delta_1<\frac12$, by ($\Gamma$2).
\item By ($\Gamma$2) at $g$ and the proof of Proposition~\ref{prop:localized-pieces}, which uses the old activities only through the bounds of Lemma~\ref{lem:gl-composed}, the pieces satisfy $|E_{\rm p}(Z)|\le\epsilon_{\rm p}e^{-(\kappa'-2)|Z|}$ and $|E_{\rm m}(Z)|\le\epsilon_{\rm m}e^{-(\kappa_n-1)d_k(Z)}$, with $\epsilon_{\rm m}=C_P\tau C_{\rm old}\mathsf K$. They depend on $\xi$ only on the bonds of $\mathfrak F$ in $Z^{(r')}$, and on $U'$ only on $Z^{(r')}$.
\item So, as in Step~3 of the proof of Proposition~\ref{prop:flow-continuity}(b), the system satisfies the hypothesis of Theorem~\ref{thm:mayer} with $\kappa=\kappa_n-1\ge C_d$ and $\epsilon=\epsilon_{\rm p}+\epsilon_{\rm m}e^{c_\natural(\kappa_n-1)}\le\epsilon_*$.
\item At $k=0$ there are only the pure pieces of the first step. They satisfy the bound of Proposition~\ref{prop:localized-pieces} with $\epsilon^{(0)}_{\rm p}\le\frac12\epsilon_*$, by ($\Gamma$2) for the first step and the end of the proof of Proposition~\ref{prop:first-step}, and the hypotheses of Theorem~\ref{thm:mayer} hold with the constants of the first step.
\item The restrictions $B_{\tilde X}$, $B_{\tilde c}$ lie in $\mathcal U_{\rm fl}(\tilde X)$ and in the corresponding set for $\tilde c$, by Remark~\ref{rem:domain-constants}(i) and (DC1), as in the first paragraph of the proof of Lemma~\ref{lem:raw-new}.
\end{itemize}
If $\hat Y=Y'$, then $Y^{(r')}\subset Y'$ (Remark~\ref{rem:extension}, collar inclusions). By the formula for $E^\#(Y)$ in the proof of Theorem~\ref{thm:mayer}, the output at $Y$ involves only the pieces $E(Z)$ with $Z\subset Y$, and these have $Z^{(r')}\subset Y'$. The restriction to the polymers with $Z^{(r')}\subset Y'$ keeps every datum of the system inside $Y'$. Its lift has diameter less than half the side.

\begin{proposition}[The step on the infinite lattice]\label{prop:gl-functional}
In the setting of Definition~\ref{def:gl-functional}:
\begin{enumerate}[label=(\alph*)]
\item \emph{(Torus independence.)} $\mathscr S_{k,\tilde Y'}(g;\mathcal E^{\rm fl};B)$ does not depend on the admissible side $n$.
\item \emph{(Covariance.)} Let $a\in M\Z^4$ in the units of scale $k+1$, and let $a_k=La$ be the same vector in the units of step $k$. Then
\[
\mathscr S_{k,\tilde Y'+a}(g;\tau_{a_k}\mathcal E^{\rm fl};\tau_aB)=\mathscr S_{k,\tilde Y'}(g;\mathcal E^{\rm fl};B).
\]
\item \emph{(The raw generation.)} Let $\tilde n\in\mathfrak N_{k+1}$ with $\operatorname{diam}\tilde Y'\le\frac12\tilde n-8$, and consider the step from scale $k$ to scale $k+1$ on the torus whose lattice of scale $k+1$ has side $\tilde n$, at the coupling $g$. If $k\ge1$, let $\mathcal E_k$ be old activities of step $k$ on it with the hypotheses of Lemma~\ref{lem:raw-new} and $\|\mathcal E_k\|\le C_{\rm old}\mathsf K$; for instance those of Theorem~\ref{thm:gen-step}(a). Let $E'=E'^{(k+1)}[\mathcal E_k,g]$ be the raw new generation (Definition~\ref{def:raw-new}). Then, with $\mathcal E^{\rm fl}_{\tilde Y'}$ the flat restriction of $\mathcal E_k$,
\[
E'(\pi_{\tilde n}\tilde Y',V_B[\tilde n])=\mathscr S_{k,\tilde Y'}(g;\mathcal E^{\rm fl}_{\tilde Y'};B)\qquad\text{for every }B\in\mathcal U_1(\tilde Y').
\]
\item \emph{(Grid covariance.)} In (c), if $k\ge1$, suppose that $\mathcal E_k(Y+a,\tau_aU)=\mathcal E_k(Y,U)$ for every polymer $Y$ of scale $k$ that does not wrap, every $a\in M\Z^4$ in the units of scale $k$, and every $U\in\mathcal U^\flat(Y)$. Then $E'(Y'+a,\tau_aV)=E'(Y',V)$ for every polymer $Y'$ of scale $k+1$ that does not wrap, every $a\in M\Z^4$ in the units of scale $k+1$, and every $V\in\mathcal U_1(Y')$. The old activities of Definition~\ref{def:old-activities} built from generations with (G2) have this covariance.
\end{enumerate}
\end{proposition}

\begin{proof}\hasdetail{prop:gl-functional}
(a) Let $n$, $n'$ be admissible for $\tilde Y'$, and let $\iota$ be the transfer of the lift $\tilde Y'$ with the vector $0$, a grid transfer with $P=Y'$ (Section~\ref{subsec:gl-setting}).
\begin{itemize}
\item \emph{References.} By Lemma~\ref{lem:gl-flat}(b), both references equal $V_B$ on the bonds of $P$, read in the lift. So they agree on $P$.
\item \emph{Pieces.} The flat data $\mathcal E_X$ of Definition~\ref{def:gl-functional} on the two tori are values of the same $\mathcal E^{\rm fl}$ at the lifts, so they agree under $\iota$. The map $Z\mapsto\iota Z$ is a bijection of the polymers of step $k$ in $P$ onto those in $\iota(P)$. It preserves the numbers of boxes and the touching relation, and $Z^{(r')}\subset Y'$ if and only if $(\iota Z)^{(r')}\subset\iota(Y')$, since $r'<R_{\rm g}$ (Lemma~\ref{lem:grid-transfer}(a), Lemma~\ref{lem:local-inventory}(a)). By Lemma~\ref{lem:gl-pieces}(a),(b), $E'(\iota Z)(\iota_*\xi)=E(Z)(\xi)$ for every such $Z$. The bonds of $\mathfrak F$ in $P$ are mapped onto those of $\mathfrak F'$ in $\iota(P)$ (Lemma~\ref{lem:local-inventory}(a)).
\item \emph{Outputs.} By the proof of Theorem~\ref{thm:mayer},
\[
E^\#(Y)=\sum_{\mathcal C:\,\bigcup_{\Gamma\in\mathcal C}\operatorname{supp}\Gamma=Y}\phi^T(\mathcal C)\prod_{\Gamma\in\mathcal C}k(\Gamma),\qquad k(\Gamma)=\int\prod_{Z\in\Gamma}K(Z)\,d\mu^*,
\]
with $K(Z)=e^{E(Z)}-1$. The clusters $\mathcal C$ consist of objects $\Gamma$, finite sets of polymers connected for touching, with supports in $Y$. So for $Y\subset P$ the bijection above induces bijections of the objects and of the clusters, with the same incompatibility graphs and hence the same Ursell coefficients. The integrand of $k(\Gamma)$ depends on $\xi$ only on the bonds of $\mathfrak F$ in $\bigcup_{Z\in\Gamma}Z^{(r')}\subset P$, and $d\mu^*$ is the product over these bonds of the standard Gaussian measures restricted to $|\xi(b)|\le p_0(g)$ and normalized. So $k'(\iota\Gamma)=k(\Gamma)$. The series converge absolutely, by the Koteck\'y--Preiss bound in the proof of Theorem~\ref{thm:mayer}. Hence $E'^\#(\iota Y;U'')=E^\#(Y;U')$, and in the same way at the reference $1$.
\item \emph{Index sets.} The rule $Y\mapsto\hat Y$ is local (Remark~\ref{rem:reblock-margins}), so $\iota$ maps the polymers $Y$ with $\hat Y=Y'$ onto those with $\hat Y=\iota(Y')$. The first two sums in Definition~\ref{def:gl-functional} are defined on the infinite lattice.
\end{itemize}

(b) The same argument applies to the transfer of the lift $\tilde Y'$ with the vector $a_k$, from $\pi_n\tilde Y'$ onto $\pi_n(\tilde Y'+a_k)$ on the same torus $\Tf[n]$. The reference of $\tau_aB$ is $\tau_{a_k}$ of that of $B$, and the flat data $\tau_{a_k}\mathcal E^{\rm fl}$ agree with $\mathcal E^{\rm fl}$ under this transfer. For the one-loop part, Proposition~\ref{prop:one-loop-local}(d) gives the covariance of $E^\infty_k$ and $J^\infty_k$, and the closure commutes with the translation, as in the proof of Lemma~\ref{lem:OL}.

(c) By Definition~\ref{def:raw-new} and Lemma~\ref{lem:raw-new}, with $V=V_B[\tilde n]$,
\[
E'(Y',V)=\sum_{\bar X=Y'}\mathfrak e(X,V)+\sum_{Y'_c=Y'}\log\mathfrak j_c(V)+\sum_{\hat Y=Y'}\big(E^\#(Y,V)-E^\#(Y,1)\big),
\]
where the outputs $E^\#(Y,V)$ are those of the restricted system of Remark~\ref{rem:extension} for $Y'$, at an extension of $V$. They do not depend on the extension (Remark~\ref{rem:extension}). We take the representation $(1,B)$ and its cutoff extension $U'$.
\begin{itemize}
\item \emph{One-loop part.} The polymers $X$ with $\bar X=Y'$ lie in $Y'$. By Remark~\ref{rem:domain-constants}(i) and (DC1), $V$ agrees on $X$ with $V_{B_{\tilde X}}$, where $B_{\tilde X}\in\mathcal U_{\rm fl}(\tilde X)$. So $\mathfrak e(X,V)=E^\infty_k(\tilde X,B_{\tilde X})$ by Proposition~\ref{prop:one-loop-local}(d), and likewise for the Jacobian terms. Lifting is a bijection of the index sets, as in the proof of Lemma~\ref{lem:OL}.
\item \emph{Outputs.} The restricted system for $Y'$ consists of the pieces $E(Z)$, $Z\subset Y'$. By the remark after Definition~\ref{def:gl-functional}, the outputs with $\hat Y=Y'$ involve only the pieces with $Z^{(r')}\subset Y'$. These depend on the reference only on $Z^{(r')}$ (Propositions~\ref{prop:interaction-pieces}(d), \ref{prop:gf-pieces}(d), \ref{prop:composed}(b), Lemma~\ref{lem:localization}(a); at $k=0$ the factorization of the composed pieces in the proof of Proposition~\ref{prop:first-step} replaces Proposition~\ref{prop:interaction-pieces}(d)), where $U'$ equals $V$. Their pure parts are those of Definition~\ref{def:gl-functional} on the torus of side $L\tilde n$ at step $k$. Their mixed parts are built from $f(X;s,Z)=\mathcal E_k(X,e^{\ii\eta A_*(s;\hat Z(Z))}U')-\mathcal E_k(X,U')$ (Proposition~\ref{prop:composed}). In the proof of that proposition and in Lemma~\ref{lem:localization}, $\mathcal E_k(X,\cdot)$ is evaluated only at configurations with real part $1$, that is, at points of $\mathcal U^\flat_1(X)$ (proof of Lemma~\ref{lem:gl-composed}). There it equals $\mathcal E^{\rm fl}_{\tilde Y'}(\tilde X,\cdot)$ at the lift. So the two systems have the same pieces with $Z^{(r')}\subset Y'$, and the outputs with $\hat Y=Y'$ coincide.
\item \emph{Normalization.} At the real critical reference $1$, the restricted and the full systems have the same outputs (Remark~\ref{rem:extension}), so $E^\#(Y,1)$ is the output of Definition~\ref{def:gl-functional} at $B=0$.
\end{itemize}
Hence $E'(Y',V)$ is the value of Definition~\ref{def:gl-functional} computed with $n=L\tilde n$, which is admissible for $\tilde Y'$. By (a), this value is $\mathscr S_{k,\tilde Y'}(g;\mathcal E^{\rm fl}_{\tilde Y'};B)$.

(d) Let $Y'$ not wrap, let $\tilde Y'$ be a lift, $V=V_B[\tilde n]$, and $a\in M\Z^4$ in the units of scale $k+1$, with $a_k=La$. Then $Y'+a$ does not wrap, $\tilde Y'+a$ is a lift of it, and $\tau_aV=V_{\tau_aB}[\tilde n]$. For $\tilde Y\subset\tilde Y'$, the hypothesis gives
\[
\mathcal E_k(\pi\tilde Y+a_k,\tau_{a_k}V_{B''})=\mathcal E_k(\pi\tilde Y,V_{B''}),
\]
since $a_k\in M\Z^4$ in the units of scale $k$. So the flat restriction of $\mathcal E_k$ to $\tilde Y'+a$ is $\tau_{a_k}\mathcal E^{\rm fl}_{\tilde Y'}$. By (c) and (b), $E'(Y'+a,\tau_aV)=E'(Y',V)$.

For the last statement, let $a\in M\Z^4$ in the units of scale $k$, and let $Y$ not wrap. In the units of scale $l\le k$, $a\in L^{k-l}M\Z^4\subset M\Z^4$. So (G2) gives $E^{(l)}(X+a,\tau_aU)=E^{(l)}(X,U)$ for $U\in\mathcal U^\flat(X)$.
\begin{itemize}
\item The translation by $a$ maps the grid cubes of every scale $l\le k$ onto grid cubes and $\Box^+$ onto $(\Box+a)^+$. It preserves $d_l$, the classes of Definition~\ref{def:polymer-classes} and the polymers that do not wrap, and $\operatorname{cl}_k(X+a)=\operatorname{cl}_kX+a$.
\item For the polymers that do not wrap, the base cubes satisfy $Q(X+a)=Q(X)+a$ (Section~\ref{sec:flow}, base cubes). The polymers $X\subset Y$ do not wrap, and the small polymers do not wrap (Lemma~\ref{lem:small-polymers}(a)).
\item $S_{\Box+a}(\tau_aU)=S_\Box(U)$ (Section~\ref{sec:flow}), and $\tau_a$ maps $\mathcal U^\flat(Y)$ onto $\mathcal U^\flat(Y+a)$ (Section~\ref{sec:generations}).
\end{itemize}
So $\Psi^{(l)}_{\Box+a}(\tau_aU)=\Psi^{(l)}_\Box(U)$ and $E^{(l),\rm ne}_{Y+a}(\tau_aU)=E^{(l),\rm ne}_Y(U)$. Hence $\mathcal E_k(Y+a,\tau_aU)=\mathcal E_k(Y,U)$.
\end{proof}

\begin{lemma}[Continuity in the old data]\label{lem:gl-continuity}
Let $k\ge1$ and $g\in(0,\gamma]$. Let $\tilde Y'$ be a polymer of scale $k+1$ of the infinite lattice, and let $\mathcal E^{\rm fl}_l\in\mathfrak D_k(\tilde Y')$ for $l\in\N\cup\{\infty\}$. Suppose that $\mathcal E^{\rm fl}_l(\tilde Y,B'')\to\mathcal E^{\rm fl}_\infty(\tilde Y,B'')$ as $l\to\infty$, for every $\tilde Y\subset\tilde Y'$ and every $B''\in\mathcal U^\flat_1(\tilde Y)$. Then
\[
\mathscr S_{k,\tilde Y'}(g;\mathcal E^{\rm fl}_l;B)\to\mathscr S_{k,\tilde Y'}(g;\mathcal E^{\rm fl}_\infty;B)\qquad\text{for every }B\in\mathcal U_1(\tilde Y').
\]
\end{lemma}

\begin{proof}\hasdetail{lem:gl-continuity}
Fix an admissible side $n$ and compute every $\mathscr S$ on $\Tf[n]$ (Proposition~\ref{prop:gl-functional}(a)). Only the mixed pieces depend on $l$. The reference $U'$, the fields $Z(s)$ and $\hat Z(Z(s))$, and $A_*(s;\hat Z(Z(s)))$ do not. We follow Steps~2--4 of the proof of Proposition~\ref{prop:flow-continuity}(b), with the couplings fixed.

\emph{Pieces.} Fix $s\in\mathbb D_{\kappa'}$, $\xi$ with $|\xi|_\infty\le p_0(g)$, and an old polymer $X\subset Z$. Then $f_l(X;s,Z(s))=\mathcal E^{\rm fl}_l(\tilde X,\tilde B_1)-\mathcal E^{\rm fl}_l(\tilde X,\tilde B_0)$, where $\tilde B_1$ and $\tilde B_0$ are the lifts of the flat fields, on $X^{(3)}$, of the composed configuration and of $U'$. They lie in $\mathcal U^\flat_1(\tilde X)$ by Lemma~\ref{lem:gl-composed}(a), and they do not depend on $l$. So $f_l(X;s,Z(s))\to f_\infty(X;s,Z(s))$, and $|f_l|\le4\tau C_{\rm old}\mathsf K$. As in Step~2 of that proof, Cauchy's formula on the circles of radius $1$ about the coordinates of the points $s(Z,t)$ writes $f_{l,\rm loc}(X,Z)$ as an integral of $f_l$ over a compact set, against a bounded kernel that does not depend on $l$. By dominated convergence, $f_{l,\rm loc}(X,Z)\to f_{\infty,\rm loc}(X,Z)$, so the mixed pieces converge, for every $\xi$ with $|\xi|_\infty\le p_0(g)$.

\emph{Outputs.} The number $k_l(\Gamma)=\int\prod_{Z\in\Gamma}K_l(Z)\,d\mu^*$ is an integral over the finitely many bonds of $\mathfrak F$ in $\bigcup_{Z\in\Gamma}Z^{(r')}$. Its integrand converges at every $\xi$ in the support of $d\mu^*$, and it is bounded by $\prod_{Z\in\Gamma}2\epsilon_*e^{-\kappa(d_k(Z)+c_\natural)}$, uniformly in $l$. By dominated convergence, $k_l(\Gamma)\to k_\infty(\Gamma)$. The proof of Theorem~\ref{thm:mayer} bounds $|k_l(\Gamma)|$ by $\bar k(\Gamma)=\prod_{Z\in\Gamma}2\epsilon_*e^{-\kappa(d_k(Z)+c_\natural)}$, which does not depend on $l$, and the Koteck\'y--Preiss theorem gives $\sum_{\mathcal C}|\phi^T(\mathcal C)|\prod_{\Gamma\in\mathcal C}\bar k(\Gamma)<\infty$, the sum over the clusters with union $Y$. By the Weierstrass test, $E^\#_l(Y;U')\to E^\#_\infty(Y;U')$, and in the same way at the reference $1$. The sum in Definition~\ref{def:gl-functional} is finite.
\end{proof}

\subsection{Convergence over the scales}\label{subsec:gl-induction}

\begin{definition}[Convergence at flat points]\label{def:gl-convergent}
Let $i\ge1$, and let $\mathcal U_*$ be $\mathcal U_1$ or $\mathcal U^\flat_1$. For $m\in\mathfrak N_i$ let $E[m]$ be a family indexed by the polymers of scale $i$ of the torus whose lattice of scale $i$ has side $m$. Suppose that for every such $X$ that does not wrap, $E[m](X,\cdot)$ is a holomorphic function on $\mathcal U_*(X)$ of the flat field that depends only on the bonds of $X$, and that
\[
\mathsf K'=\sup_m\ \sup_X\ e^{\kappa_nd_i(X)}\sup_{\mathcal U_*(X)}|E[m](X,\cdot)|<\infty,
\]
the second supremum over the polymers that do not wrap. The collection $(E[m])_{m\in\mathfrak N_i}$ \emph{converges on $\mathcal U_*$} if, for every polymer $\tilde X$ of scale $i$ of the infinite lattice and every $B\in\mathcal U_*(\tilde X)$, the numbers $E[m](\pi_m\tilde X,V_B[m])$ converge as $m\to\infty$ in $\mathfrak N_i$. These numbers are defined for the $m$ with $\operatorname{diam}\tilde X\le\frac12m-8$. The limit $E^\infty(\tilde X,B)$ is the \emph{limit} of the collection.
\end{definition}

This is a qualitative counterpart of \eqref{eq:torus-local}, on $\mathcal U_1$ or $\mathcal U^\flat_1$; \eqref{eq:torus-local} itself is not changed.

\begin{lemma}[Limits of convergent families]\label{lem:gl-limits}
Let $(E[m])_{m\in\mathfrak N_i}$ converge on $\mathcal U_*$, with limit $E^\infty$.
\begin{enumerate}[label=(\alph*)]
\item Each $E^\infty(\tilde X,\cdot)$ is holomorphic on $\mathcal U_*(\tilde X)$, depends only on $B$ on the bonds of $\tilde X$, and satisfies $|E^\infty(\tilde X,\cdot)|\le\mathsf K'e^{-\kappa_nd_i(\tilde X)}$. The convergence is uniform on the compact subsets of $\mathcal U_*(\tilde X)$. If the functions $E[m](X,\cdot)$ vanish at $0$, are real for real $B$, or are invariant under small real gauge transformations in the sense of Section~\ref{subsec:ol-period-vacpol}, then so is $E^\infty(\tilde X,\cdot)$.
\item If $E[m](X+a,\tau_aV)=E[m](X,V)$ for every $m$, every polymer $X$ that does not wrap, every $a\in M\Z^4$ and every $V\in\mathcal U_*(X)$, then $E^\infty(\tilde X+a,\tau_aB)=E^\infty(\tilde X,B)$.
\item Suppose, in addition to the hypothesis of (b), that every $E[m]$ is a family of normalized graded activities at the polymers of scale $i$, on domains containing the $\mathcal U^\flat(X)$, real at real configurations and invariant under real gauge transformations, with $\mathsf K''=\sup_m\|E[m]\|<\infty$ (the norm of Section~\ref{subsec:ol-flow-coeff}, over all polymers). Then:
  \begin{itemize}
  \item the restriction of $E^\infty$ to the domains $\mathcal U^\flat_1$ is a family of flat activities with (G2), and $\|E^\infty\|\le\mathsf K''$;
  \item $\beta^{\rm vol}_{E[m]}\to\beta^\infty_{E^\infty}$ as $m\to\infty$ in $\mathfrak N_i$;
  \item the kernels of the second variations at $A=0$ of $A\mapsto\sum_XE[m](X,e^{\ii\eta A})$ converge to the kernel $\widetilde\Pi$ of $E^\infty$ (Section~\ref{subsec:ol-period-vacpol}) with uniform decay, in the sense of Lemma~\ref{lem:ol-folding}(e).
  \end{itemize}
\end{enumerate}
\end{lemma}

\begin{proof}\hasdetail{lem:gl-limits}
(a) Let $\operatorname{diam}\tilde X\le\frac12m-8$. Then $d_i(\pi_m\tilde X)=d_i(\tilde X)$.
\begin{itemize}
\item Projecting a tree of $\R^4$ that meets the cubes of $\tilde X$ gives a tree of the same length that meets those of $\pi_m\tilde X$. So $d_i(\pi_m\tilde X)\le d_i(\tilde X)$.
\item Conversely, the projections to the coordinate circles of the union of the closed cubes of $\pi_m\tilde X$ are arcs (Section~\ref{sec:flow}, base cubes), of length at most $\operatorname{diam}\tilde X+\eta<\frac12m$. Map each circle onto its arc by the identity on the arc, and on the complement, which is longer than the arc, by a $1$-Lipschitz map onto the arc that fixes the two endpoints. This map is $1$-Lipschitz. The product of these maps is a $1$-Lipschitz retraction of the torus onto a box that contains the cubes and lifts isometrically.
\item It maps a tree that meets the cubes onto a connected finite union of rectifiable curves, of no greater length, that meets the same cubes and contains a tree that meets them. So $d_i(\tilde X)\le d_i(\pi_m\tilde X)$.
\end{itemize}
The functions $B\mapsto E[m](\pi_m\tilde X,V_B[m])$ are then holomorphic on the open connected set $\mathcal U_*(\tilde X)$ and bounded by $\mathsf K'e^{-\kappa_nd_i(\tilde X)}$. By Montel's theorem every subsequence has a further subsequence that converges uniformly on compact subsets to a holomorphic function, and this function is the pointwise limit $E^\infty(\tilde X,\cdot)$. So the whole sequence converges uniformly on compact subsets, and $E^\infty(\tilde X,\cdot)$ is holomorphic. The bound, the dependence on the bonds of $\tilde X$, the vanishing at $0$, the reality and the invariance pass to the pointwise limit.

(b) $\pi_m(\tilde X+a)=\pi_m\tilde X+a$ and $V_{\tau_aB}[m]=\tau_aV_B[m]$.

(c) \emph{Flat activities.} The torus functions vanish at $0$ and are real for real $B$. They are invariant under small real gauge transformations, as in the proof of Proposition~\ref{prop:one-loop-local}(d): extend the gauge function from $\pi_m\tilde X$ to the torus, and use the invariance of $E[m](X,\cdot)$ under real gauge transformations and its dependence on the bonds of $X$ only. By (a) and (b), the restriction of $E^\infty$ to $\mathcal U^\flat_1$ is a family of flat activities with (G2). Moreover $\|E^\infty\|\le\min(\mathsf K',\mathsf K'')$, since $\mathcal U^\flat_1(X)\subset\mathcal U^\flat(X)\cap\mathcal U_*(X)$ and $d_i(\pi_m\tilde X)=d_i(\tilde X)$ by (a).

\emph{The coefficient.} For a polymer $X$ that does not wrap and small $|t|$, the configurations $V^X_{t,F}$ of Lemma~\ref{lem:period-density}(a) are points of $\mathcal U^\flat_1(X)$, and $\tau_aV^X_{t,F}=V^{X+a}_{t,F}$ (proof of Lemma~\ref{lem:period-density}(b)). So the hypothesis of (b) gives $c^{X+a}_{E[m]}=c^X_{E[m]}$. The sums $\sum_{X\supset Q}c^X_{E[m]}(F)/|X|$ over the polymers that do not wrap are then the same for all grid cubes $Q$, and the paragraph ``The volume coefficient'' of Section~\ref{subsec:ol-flow-coeff} gives
\[
\bar T^{\rm vol}_{E[m]}(F)=M^{-4}\sum_{X\supset Q}\frac{c^X_{E[m]}(F)}{|X|}
\]
for one fixed cube $Q$. Lifting with $Q\mapsto\tilde Q$ is a bijection from the polymers $X\supset Q$ that do not wrap onto the polymers $\tilde X\supset\tilde Q$ of the infinite lattice of diameter at most $\frac12m-8$ (proof of Lemma~\ref{lem:period-vacpol}(d)). For fixed $\tilde X$ and $m$ large:
\begin{itemize}
\item $c^{\pi_m\tilde X}_{E[m]}(F)$ is a Cauchy integral over a fixed circle in $t$, as in Step~5 of the proof of Proposition~\ref{prop:flow-continuity}(b). By the uniform convergence on compact sets in (a), it converges to $c^{\tilde X}_{E^\infty}(F)$.
\item By Lemma~\ref{lem:period-density}(a), the terms are bounded by $9d^2M^2(1+d_i(\tilde X))^2(\delta_1^\flat)^{-2}\mathsf K''e^{-\kappa_nd_i(\tilde X)}|F|^2$. This is summable over $\tilde X\supset\tilde Q$ by \eqref{eq:dM-facts}, since $(1+t)^2e^{-(\kappa_n-\kappa_0)t}\le e^{-(\kappa_n-\kappa_0-2)t}$ and $\kappa_n-\kappa_0\ge3$.
\end{itemize}
By dominated convergence, $\bar T^{\rm vol}_{E[m]}(F)\to\bar T_{E^\infty}(F)$ for every $F$, and so $\beta^{\rm vol}_{E[m]}\to\beta^\infty_{E^\infty}$ by \eqref{eq:beta-vol} and the definition of $\beta^\infty_{E^\infty}$.

\emph{The kernels.} Fix two bonds of the infinite lattice and their projections $b,b'$. The torus kernel at $(b,b')$ is the sum, over the polymers $X$ of the torus that contain both, of the kernels of the second variations of $E[m](X,e^{\ii\eta A})$. Only small fields $A$ enter, so only points of $\mathcal U^\flat_1(X)$ for the polymers that do not wrap. These are:
\begin{itemize}
\item \emph{Polymers that do not wrap.} Their kernels are $2\kappa_X(b,b')$, in the notation of Lemma~\ref{lem:period-vacpol}(a), and they correspond to the polymers $\tilde X$ of the infinite lattice of diameter at most $\frac12m-8$ that contain the lifts. For fixed $\tilde X$ and $m$ large, $\kappa_{\pi_m\tilde X}\to\kappa_{\tilde X}$, by Cauchy's estimate and the uniform convergence on compact sets.
\item \emph{All polymers.} By the same Cauchy estimate, each kernel is bounded by $C_\eta\mathsf K''e^{-\kappa_nd_i(X)}$, and this is summable by \eqref{eq:dM-facts}.
\item \emph{Wrapping polymers.} They have $d_i(X)\ge m/(4M)$ (Lemma~\ref{lem:torus}(b)), so they contribute at most $C_\eta\mathsf K''K_0e^{-(\kappa_n-\kappa_0)m/(4M)}\to0$.
\end{itemize}
By dominated convergence, the torus kernels converge entrywise to $\widetilde\Pi=2\sum_{\tilde X}\kappa_{\tilde X}$. The bound $C_\eta K_0\mathsf K''e^{2(\kappa_n-\kappa_0)}e^{-(\kappa_n-\kappa_0)|m_b-m_{b'}|_m/M}$, uniform in $m$, is proved as at the end of the proof of Lemma~\ref{lem:period-vacpol}(d), and it uses only $\mathsf K''$.
\end{proof}

\begin{theorem}[Generation locality]\label{thm:gen-locality}
Let $G$ be compact, connected and simple, let $\gamma$ satisfy ($\Gamma$1)--($\Gamma$3), let $j\ge1$ and $\pi=(g_0,\dots,g_{j-1})\in\Sigma_{j-1}$, and use the notation of Section~\ref{subsec:gl-setting}. Then for every $i\le j$:
\begin{enumerate}[label=(\alph*)]
\item For every $m\in\mathfrak N_i$, $E'^{(i)}[m]$ has the covariance of Lemma~\ref{lem:gl-limits}(b) on $\mathcal U_1$. The collection $(E'^{(i)}[m])_{m\in\mathfrak N_i}$ converges on $\mathcal U_1$, with $\mathsf K'\le\mathsf K$. Its limit $E'^{(i),\infty}$ is covariant under the translations by $M\Z^4$, vanishes at $0$, is real for real $B$, and is invariant under small real gauge transformations.
\item The collection $(E^{(i)}[m])_{m\in\mathfrak N_i}$ converges on $\mathcal U^\flat_1$, with the limit
\[
E^{(i),\infty}(\tilde X,B)=\frac1{|B_4|}\sum_{R\in B_4}E'^{(i),\infty}\big(R^{-1}\tilde X,R^{-1}B\big),
\]
with $B_4$ acting about the origin.
\item $\beta_i[m]\to\beta^\infty_i:=\beta^\infty_{E'^{(i),\infty}}$ as $m\to\infty$ in $\mathfrak N_i$.
\item If $i<j$, then, for every polymer $\tilde Y'$ of scale $i+1$ of the infinite lattice, the flat restrictions $\mathcal E^{\rm fl}_{\tilde Y'}$ of the old activities $\mathcal E_i[m']$ converge, as $m'\to\infty$ in $\mathfrak N_{i+1}$, at every point of the domains $\mathcal U^\flat_1(\tilde Y)$, $\tilde Y\subset\tilde Y'$, to the flat old data $\mathcal E^\infty_i\in\mathfrak D_i(\tilde Y')$. These data are built by Definition~\ref{def:old-activities} on the infinite lattice from the $E^{(l),\infty}$ and the $\beta^\infty_l$, $l\le i$.
\end{enumerate}
The limits depend only on $(g_0,\dots,g_{i-1})$.
\end{theorem}

\begin{proof}\hasdetail{thm:gen-locality}
By induction on $i$. We show that (a) at all scales $l\le i$ gives (b) and (c) at the scale $i$, and (d) at the scale $i$ if $i<j$. Then we show (a) at $i=1$, and (a) at $i=k+1$ from (b)--(d) at all scales $\le k$.

\emph{(a) at scale $i$ gives (b).} Let $X$ be a polymer of scale $i$ of the torus of side $m$ that does not wrap, and let $U=V_B[m]\in\mathcal U^\flat_1(X)$. By Definition~\ref{def:symmetrization},
\[
E^{(i)}[m](X,U)=|\mathfrak S^{\rm grid}_i|^{-1}\sum_rE'^{(i)}[m](r^{-1}X,r^{-1}U).
\]
\begin{itemize}
\item Write $r(x)=Rx+a$ with $R\in B_4$, acting about the origin, and $a\in M\Z^4$. Then $r^{-1}X=R^{-1}X-R^{-1}a$ and $r^{-1}U=\tau_{-R^{-1}a}(R^{-1}U)$, with $R^{-1}a\in M\Z^4$.
\item The polymer $R^{-1}X$ does not wrap, and $R^{-1}U\in\mathcal U_1(R^{-1}X)$ (Lemma~\ref{lem:gl-flat}(c)). So the covariance in (a) gives $E'^{(i)}[m](r^{-1}X,r^{-1}U)=E'^{(i)}[m](R^{-1}X,R^{-1}U)$.
\item Each $R$ is the linear part of $(m/M)^4$ elements $r$ (proof of Proposition~\ref{prop:period-iso}(c)). So $E^{(i)}[m](X,U)=|B_4|^{-1}\sum_RE'^{(i)}[m](R^{-1}X,R^{-1}U)$.
\end{itemize}
Let $\tilde X$ be a polymer of the infinite lattice and $B\in\mathcal U^\flat_1(\tilde X)$. For $m$ large, $\pi_m\tilde X$ and the $\pi_mR^{-1}\tilde X$ do not wrap, and $R^{-1}V_B[m]=V_{R^{-1}B}[m]$ with $R^{-1}B\in\mathcal U_1(R^{-1}\tilde X)$ (Lemma~\ref{lem:gl-flat}(c)). By (a), each term converges, which gives (b). The bound $\mathsf K'\le\mathsf K$ on $\mathcal U^\flat_1$ holds since $\|E^{(i)}[m]\|_i\le\mathsf K$ (Proposition~\ref{prop:flow-continuity}(a)).

\emph{(a) at scale $i$ gives (c).} By Proposition~\ref{prop:period-iso}(c) and the paragraph ``The volume coefficient'' of Section~\ref{subsec:ol-flow-coeff}, $\beta_i[m]=\beta^{\rm vol}_{E'^{(i)}[m]}$. The raw generations are normalized graded activities on the domains $\mathcal U(X)\supset\mathcal U^\flat(X)$, real at real configurations and gauge invariant (Lemma~\ref{lem:raw-new}(a)), with $\|E'^{(i)}[m]\|\le\|E'^{(i)}[m]\|^{\rm raw}_i\le\mathsf K$ (Theorem~\ref{thm:gen-step}(b); Proposition~\ref{prop:gen-one} for $i=1$). Lemma~\ref{lem:gl-limits}(c) applies on $\mathcal U_1$, with the covariance of (a), and gives (c). Applied on $\mathcal U^\flat_1$ to the generations $E^{(i)}[m]$, which have (G2) and (G4), so that $\beta_{E^{(i)}[m]}=\beta^{\rm vol}_{E^{(i)}[m]}$ (Section~\ref{subsec:ol-flow-coeff}), it also gives $\beta^\infty_i=\beta^\infty_{E^{(i),\infty}}$.

\emph{(a)--(c) at the scales $l\le i$ give (d).} Let $\tilde Y'$ be a polymer of scale $i+1$, $\tilde Y\subset\tilde Y'$ a polymer of scale $i$, and $B''\in\mathcal U^\flat_1(\tilde Y)$. On the torus whose lattice of scale $i+1$ has side $m'$, the side at scale $l$ is $L^{i+1-l}m'\in\mathfrak N_l$. By Definition~\ref{def:old-activities},
\[
\mathcal E_i[m'](\pi\tilde Y,V_{B''})=\sum_{l=1}^i\Big(\one_{\{\pi\tilde Y=\Box^+\}}\Psi^{(l)}_\Box(V_{B''})+E^{(l),\rm ne}_{\pi\tilde Y}(V_{B''})\Big).
\]
\begin{itemize}
\item \emph{Finite sums.} The polymers $X$ in these sums lie in $\pi\tilde Y$: the small polymers based in $\Box$ lie in $\Box^+$ (Lemma~\ref{lem:small-polymers}(b)), and the others have $\operatorname{cl}_iX=\pi\tilde Y$. So the sums are finite.
\item \emph{Classes.} For $m'$ large, $\pi\tilde Y$ does not wrap. The polymers $X\subset\pi\tilde Y$ of scale $l$ then correspond, under lifting, to the polymers $\tilde X\subset\tilde Y$ of the infinite lattice, with the same classes, closures and base cubes (Section~\ref{subsec:gl-setting}, the infinite lattice).
\item \emph{Terms.} By Lemma~\ref{lem:gl-flat}(a), $V_{B''}$ is on $\tilde X^{(3)}$ a point of $\mathcal U^\flat_1(\tilde X)$ of scale $l$. So each term $E^{(l)}[L^{i+1-l}m'](X,V_{B''})$ converges by (b) at the scale $l$, along a subsequence of $\mathfrak N_l$. Each $\beta_l[L^{i+1-l}m']$ converges by (c), and $S_\Box(V_{B''})$ does not depend on $m'$.
\end{itemize}
So $\mathcal E^{\rm fl}_{\tilde Y'}(\tilde Y,B'')$ converges, and the limit is the value of the data $\mathcal E^\infty_i$ built by Definition~\ref{def:old-activities} on the infinite lattice. These data are finite sums of holomorphic functions on $\mathcal U^\flat_1(\tilde Y)$, by Lemma~\ref{lem:gl-flat}(a). They vanish at $0$, since the $E^{(l),\infty}$ do (Lemma~\ref{lem:gl-limits}(a)) and $S_\Box(1)=0$. They satisfy the bound $C_{\rm old}\mathsf Ke^{-\kappa_nd_i(\tilde Y)}$, as limits of the values of the torus activities, which satisfy it (Proposition~\ref{prop:old-activities}(b)). So $\mathcal E^\infty_i\in\mathfrak D_i(\tilde Y')$.

\emph{(a) at $i=1$.} Let $\tilde X$ be a polymer of scale $1$ of the infinite lattice and $B\in\mathcal U_1(\tilde X)$. For $m\in\mathfrak N_1$ with $\operatorname{diam}\tilde X\le\frac12m-8$, Proposition~\ref{prop:gl-functional}(c) at $k=0$ gives $E'^{(1)}[m](\pi_m\tilde X,V_B[m])=\mathscr S_{0,\tilde X}(g_0;B)$, which does not depend on $m$. The covariance on $\mathcal U_1$ is Proposition~\ref{prop:gl-functional}(d), with no old activities. The bound $\mathsf K'\le\mathsf K$ holds since $\mathcal U_1(X)\subset\mathcal U(X)$ and $\|E'^{(1)}[m]\|^{\rm raw}_1\le\mathsf K$ (Proposition~\ref{prop:gen-one}). The other properties follow from Lemma~\ref{lem:gl-limits}(a),(b), since $E'^{(1)}[m]$ has (G1) on $\mathcal U(X)$ (Lemma~\ref{lem:raw-new}(a)) and is invariant under small real gauge transformations, as in the proof of Lemma~\ref{lem:gl-limits}(c).

\emph{(a) at $i=k+1\le j$.} Assume (a)--(c) at the scales $l\le k$, so that (d) holds at $k$. Let $\tilde Y'$ be a polymer of scale $k+1$ of the infinite lattice and $B\in\mathcal U_1(\tilde Y')$. For $m\in\mathfrak N_{k+1}$ with $\operatorname{diam}\tilde Y'\le\frac12m-8$, the raw generation $E'^{(k+1)}[m]$ is that of Proposition~\ref{prop:flow-continuity}(a), built at the coupling $g_k$ with the old activities $\mathcal E_k[m]$. These satisfy the hypotheses of Lemma~\ref{lem:raw-new}, with $\|\mathcal E_k[m]\|\le C_{\rm old}\mathsf K$, by Theorem~\ref{thm:gen-step}(a), applied as in the proof of Proposition~\ref{prop:flow-continuity}(a). Proposition~\ref{prop:gl-functional}(c) gives
\[
E'^{(k+1)}[m](\pi_m\tilde Y',V_B[m])=\mathscr S_{k,\tilde Y'}\big(g_k;\mathcal E^{\rm fl}_{\tilde Y'}[\mathcal E_k[m]];B\big).
\]
By (d) at $k$ and Lemma~\ref{lem:gl-continuity}, the right side converges to $\mathscr S_{k,\tilde Y'}(g_k;\mathcal E^\infty_k;B)$. So the collection converges on $\mathcal U_1$. The bound $\mathsf K'\le\mathsf K$ holds since $\|E'^{(k+1)}[m]\|^{\rm raw}_{k+1}\le\mathsf K$ (Theorem~\ref{thm:gen-step}(b)).

The old activities $\mathcal E_k[m]$ are built from the generations $E^{(l)}[\cdot]$, which have (G2) (Lemma~\ref{lem:symmetrization}(a)). So Proposition~\ref{prop:gl-functional}(d) gives the covariance on $\mathcal U_1$. The remaining properties of the limit follow from Lemma~\ref{lem:gl-limits}(a),(b), as at $i=1$.

All the constructions use only $(g_0,\dots,g_{i-1})$.
\end{proof}

\begin{remark}[What is not proved]\label{rem:gl-rate}
\begin{enumerate}[label=(\alph*)]
\item \emph{No rate.} The convergence in Theorem~\ref{thm:gen-locality} has no rate, and it is not shown to be uniform in the couplings. A rate, such as $|\beta_j[\tilde n]-\beta^\infty_j|\le C'e^{-\tilde n/(4M)}$, is expected to follow from three ingredients:
  \begin{itemize}
  \item a Lipschitz bound of $\mathscr S_{k,\tilde Y'}$ in the flat old data, with a constant proportional to $\tau$, by an interpolation in the mixed pieces as in Corollary~\ref{cor:two-classes};
  \item a bound on the errors of the old activities;
  \item for constants uniform in $j$, a further condition on $\gamma$.
  \end{itemize}
  Uniformity in the couplings would need these constants to be uniform on $\Sigma_{j-1}$. This is not carried out here.
\item \emph{The couplings of the flow.} In the flow of Theorem~\ref{thm:sf-flow} the couplings are computed on the given torus, so they may depend on it. Theorem~\ref{thm:gen-locality} gives $\beta_j[\tilde n]\to\beta^\infty_j$ at fixed couplings. It gives no bound on $\beta_j-\beta^\infty_j$ on a given torus, and, without uniformity in the couplings, it does not show that this difference tends to $0$ along the flows on growing tori.
\item \emph{Comparison with Ba\l aban.} Ba\l aban's generation $E^{(j+1)}$ and his $\beta_{j+1}$ are functions of $g_j$ alone, and he takes the infinite-volume limit of (1.20) ``by the localized representation (1.7)'' \cite[(1.18), (1.20)--(1.22)]{BalYM4I}. The generations of this note depend on $(g_0,\dots,g_{j-1})$ through the old activities and their counterterms $\beta_i$. This is why their convergence needs the induction over the scales.
\end{enumerate}
\end{remark}
